% =============================================================================
% Patch: §9 N-point odd-moment cancellation
%   sec:n-point-odd-moment-cancellation
%
% Target file:   rh/rh_rebuild.tex
% Target lines:  1318 through 1357 (current §9 stub: bare lemma + wip)
%
% Action: replace the existing §9 with the content below.
%
% Migration content:
%   - Explicit coefficients a_j^{(N)} = (-1)^{j+1} j C(2N, N+j) / C(2N-2, N-1).
%   - Coefficient lemma (lem:n-point-coefficients): sum a_j = 1,
%     sum a_j j^{2r+1} = 0 for 0 <= r <= N-2, sum |a_j| = 2N - 1.
%   - Discrete measure mu_hat_{N,Q}(z) := sum a_j exp(j z / Q^2),
%     odd-jet vanishing at 0 up to order 2N - 3 (lem:n-point-odd-moment-cancellation).
%   - Convolution lemma (lem:n-point-convolution-parity): even bumps
%     preserve the parity vanishing.
%   - Algebraic-difference identity (lem:A-N-difference) confirming
%     A_N(t) - A_N(1/t) = (-1)^N N / C(2N-2, N-1) * t^{-N} (1-t)^{2N-1} (1+t),
%     which exhibits the (1-t)^{2N-1} factor responsible for the order
%     2N - 1 vanishing of the odd part.
%   - Scope remark: even-bump hypothesis is load-bearing; an odd
%     contamination breaks the cancellation at the leading odd order.
%
% Fixes overclaim in the prior stub: the lemma's odd-moment vanishing is
% ONLY for 0 <= r <= N - 2, not "for every r >= 0".  The first surviving
% odd derivative of mu_hat_{N,Q} is at order 2N - 1, with explicit value
% (-1)^{N+1} (2N-1) N! (N-1)! / Q^{2(2N-1)}.
%
% Closes:
%   rem:wip-n-point-odd-moment-cancellation
%
% Status:
%   \statusmig      not-started -> migrated
%   \statusreferee  not-started -> in-progress
%   \statussympy    not-started -> verified
%       [rh/sympy/sec-n-point-odd-moment-cancellation.py]
%   \statussim      not-started -> verified
%       [rh/simulation/sec-n-point-odd-moment-cancellation.py]
%   \statuslean     not-started -> n/a
%
% Verification (every claim has a failing test in the cited script):
%   - rh/sympy/sec-n-point-odd-moment-cancellation.py
%       verify_coefficient_identities (N = 1 through 8),
%       verify_mu_hat_parity (N = 1 through 5),
%       verify_convolution_parity (N = 1 through 5, even phi up to deg 6),
%       verify_A_N_difference_identity (N = 1 through 5),
%       verify_alternating_sign_pattern (N = 1 through 8).
%   - rh/simulation/sec-n-point-odd-moment-cancellation.py
%       check_coefficient_identities (N = 1 through 12, mp.dps = 100),
%       check_mu_hat_odd_vanishing (N = 1 through 7, FD via mpmath.diff),
%       check_convolution_parity (random even bumps),
%       check_ell1_growth (N up to 50),
%       check_parity_broken_for_odd_bump (scope stress),
%       check_difference_identity_numerical (independent route).
% =============================================================================

% =============================================================================
\section{\texorpdfstring{\(N\)}{N}-point odd-moment cancellation}
\label{sec:n-point-odd-moment-cancellation}
\statuspaper{LIVE}\quad
\statusmig{migrated}\quad
\statusreferee{in-progress}\quad
\statussympy[rh/sympy/sec-n-point-odd-moment-cancellation.py]{verified}\quad
\statussim[rh/simulation/sec-n-point-odd-moment-cancellation.py]{verified}\quad
\statuslean{n/a}
\par\medskip

The two-point scalar channel, and any future four-point analogue, may
produce Taylor remainder terms in the separation variable.  This section
supplies only the algebraic cancellation package used to remove low-order
odd terms: an explicit signed discrete \(N\)-point measure
\(\mu_{N,Q}\) on \(s\)-space whose low odd moments vanish by
construction.  It does not by itself prove that the remainders are below
\(Q^{-A_\toy}\), below \(Q^{-B_\z}\), or uniformly negligible in height;
those estimates belong to the two-/four-point comparison and aggregation
sections.

\subsection{Coefficients and moment identities}
\label{subsec:n-point-coefficients}

\begin{definition}[\(N\)-point coefficients]
\label{def:n-point-coefficients}
For each integer \(N\ge 1\) and \(1\le j\le N\),
\begin{equation}
\label{eq:n-point-coefficients}
        a_j^{(N)}
        :=
        (-1)^{j+1}\,
        \frac{j\,\binom{2N}{N+j}}{\binom{2N-2}{N-1}}.
\end{equation}
\end{definition}

\begin{lemma}[Moment identities for \(a_j^{(N)}\)]
\label{lem:n-point-coefficients}
The coefficients \(a_j^{(N)}\) of \eqref{eq:n-point-coefficients} satisfy
\begin{align}
\label{eq:n-point-normalization}
        \sum_{j=1}^N a_j^{(N)}
        &= 1,\\
\label{eq:n-point-odd-moments-vanish}
        \sum_{j=1}^N a_j^{(N)}\, j^{\,2r+1}
        &= 0
        \qquad (0\le r\le N-2),\\
\label{eq:n-point-ell1}
        \sum_{j=1}^N |a_j^{(N)}|
        &= 2N - 1.
\end{align}
Consequently, for every integer \(k\ge 0\),
\begin{equation}
\label{eq:n-point-moment-bound}
        \Bigl|\sum_{j=1}^N a_j^{(N)}\, j^{\,k}\Bigr|
        \;\le\;
        (2N-1)\,N^{\,k}.
\end{equation}
\end{lemma}

\begin{proof}
Write \(D_N=\binom{2N-2}{N-1}\), so
\(a_j^{(N)} = (-1)^{j+1} j \binom{2N}{N+j} / D_N\).

\emph{Normalization.}  The identity
\[
        j\binom{2N}{N+j}
        =N\left(\binom{2N-1}{N+j-1}-\binom{2N-1}{N+j}\right)
\]
gives
\[
\sum_{j=1}^N(-1)^{j+1}j\binom{2N}{N+j}
=N T_N,
\]
where
\[
T_N:=\sum_{j=1}^N(-1)^{j+1}
\left(\binom{2N-1}{N+j-1}-\binom{2N-1}{N+j}\right).
\]
Putting \(k=N+j-1\) in the first sum and \(k=N+j\) in the second,
and using \(\binom{2N-1}{2N}=0\), one obtains
\[
        T_N=2\sum_{k=N}^{2N-1}(-1)^{k-N}\binom{2N-1}{k}
             -\binom{2N-1}{N}.
\]
The partial alternating-binomial identity gives
\[
        \sum_{k=N}^{2N-1}(-1)^{k-N}\binom{2N-1}{k}=D_N,
\]
and
\(\binom{2N-1}{N}=((2N-1)/N)D_N\).  Hence \(T_N=D_N/N\), and so
\(\sum_j a_j^{(N)}=D_N^{-1}NT_N=1\).

\emph{Odd-moment vanishing.}  For any polynomial \(P(k)\) of degree
\(<2N\),
\begin{equation}
\label{eq:n-point-finite-difference-identity}
        \sum_{k=0}^{2N}(-1)^k\binom{2N}{k}P(k) = 0,
\end{equation}
because the operator on the left is the \(2N\)-th finite difference of
\(P\) at the origin and any polynomial of degree \(<2N\) is annihilated
by it.  Apply \eqref{eq:n-point-finite-difference-identity} to
\(P(k)=(k-N)^{2r+2}\) for \(0\le r\le N-2\) (so \(2r+2 < 2N\)).  After
the change of variables \(j=k-N\), this becomes
\[
        \sum_{j=-N}^N(-1)^j\binom{2N}{N+j}j^{\,2r+2}=0.
\]
The summand is even in \(j\): the power \(j^{2r+2}\) is even,
\((-1)^{-j}=(-1)^j\), and
\(\binom{2N}{N-j}=\binom{2N}{N+j}\).  The \(j=0\) term vanishes, so
\[
        2\sum_{j=1}^N(-1)^j\binom{2N}{N+j}j^{\,2r+2}=0.
\]
Since
\[
        \sum_{j=1}^N a_j^{(N)}j^{2r+1}
        =-D_N^{-1}
          \sum_{j=1}^N(-1)^j\binom{2N}{N+j}j^{2r+2},
\]
we obtain \eqref{eq:n-point-odd-moments-vanish}.

\emph{\(\ell^1\) control.}  The signs of \(a_j^{(N)}\) alternate, so
\(\sum_{j=1}^N|a_j^{(N)}| = D_N^{-1}\sum_{j=1}^N j\binom{2N}{N+j}\).
Telescoping the same identity used for normalization gives
\(\sum_{j=1}^N j\binom{2N}{N+j}=N\binom{2N-1}{N}\), and the standard
identity \(N\binom{2N-1}{N}=(2N-1)\binom{2N-2}{N-1}\) yields
\(\sum_{j=1}^N|a_j^{(N)}|=2N-1\).

The bound \eqref{eq:n-point-moment-bound} follows from
\(|\sum a_j j^k| \le \sum|a_j|\,j^k \le (2N-1)N^k\).
\end{proof}

\begin{remark}[The case \(N=1\)]
\label{rem:n-point-N-equals-one}
For \(N=1\), the range \(0\le r\le N-2\) is empty, so there is no
low odd moment to cancel.  The normalization gives \(a_1^{(1)}=1\), and
Corollary~\ref{cor:n-point-first-surviving} says that the first odd
jet is already the first derivative.  Thus genuine cancellation starts
at \(N\ge2\), while the formulas remain valid for \(N=1\).
\end{remark}


\subsection{The discrete \(N\)-point measure}
\label{subsec:n-point-measure}

\begin{definition}[\(N\)-point measure and exponential symbol]
\label{def:n-point-measure}
For each height parameter \(Q>0\), let \(\mu_{N,Q}\) be the signed
discrete measure on \(\mathbb R\) supported on
\(\{j/Q^2:1\le j\le N\}\) with masses \(a_j^{(N)}\) of
Definition~\ref{def:n-point-coefficients}.  Its exponential symbol is
\begin{equation}
\label{eq:mu-hat-N-Q}
        \widehat\mu_{N,Q}(z)
        := \sum_{j=1}^N a_j^{(N)}\, e^{j z / Q^2}.
\end{equation}
Here \(z\) is a formal or complex local variable.  No factor of \(i\) is
included in this convention; the hat denotes the generating symbol used
for moment cancellation, not a unitary Fourier transform.
\end{definition}

\begin{lemma}[Odd-jet vanishing of \(\widehat\mu_{N,Q}\)]
\label{lem:n-point-odd-moment-cancellation}
For every \(N\ge 1\) and every \(Q>0\),
\begin{equation}
\label{eq:mu-hat-normalization}
        \widehat\mu_{N,Q}(0) = 1,
\end{equation}
and
\begin{equation}
\label{eq:mu-hat-odd-vanishing}
        \partial_z^{\,2r+1}\widehat\mu_{N,Q}(0) = 0
        \qquad (0 \le r \le N - 2).
\end{equation}
\end{lemma}

\begin{proof}
Differentiating \eqref{eq:mu-hat-N-Q},
\[
        \partial_z^{\,k}\widehat\mu_{N,Q}(0)
        = Q^{-2k}\sum_{j=1}^N a_j^{(N)}\, j^{\,k}.
\]
For \(k=0\) this is \(\sum_j a_j^{(N)} = 1\) by
\eqref{eq:n-point-normalization}.  For \(k=2r+1\) with \(0\le r\le N-2\)
the inner sum is zero by \eqref{eq:n-point-odd-moments-vanish}.
\end{proof}

\begin{lemma}[Convolution stability]
\label{lem:n-point-convolution-parity}
Let \(\widehat\phi:\mathbb R\to\mathbb R\) be even and \(C^\infty\) near
\(0\), and let \(\eta\in\mathbb R\).  Then for every \(N\ge 1\) and
every \(Q>0\),
\begin{equation}
\label{eq:n-point-convolution-vanishing}
        \partial_z^{\,2r+1}
        \bigl(\widehat\phi(\eta z)\,\widehat\mu_{N,Q}(z)\bigr)
        \big|_{z=0}
        = 0
        \qquad (0 \le r \le N - 2).
\end{equation}
\end{lemma}

\begin{proof}
By the Leibniz rule,
\[
        \partial_z^{\,2r+1}
        \bigl(\widehat\phi(\eta z)\widehat\mu_{N,Q}(z)\bigr)\big|_{z=0}
        = \sum_{a+b = 2r+1}\binom{2r+1}{a}\,
          \eta^a\,\widehat\phi^{\,(a)}(0)\,
          \partial_z^{\,b}\widehat\mu_{N,Q}(0).
\]
With \(a+b\) odd, exactly one of \(a, b\) is odd.  If \(a\) is odd,
\(\widehat\phi^{(a)}(0)=0\) since \(\widehat\phi\) is even.  If \(b\) is
odd, then \(b\le 2r+1\le 2N-3\), so
\(\partial_z^b\widehat\mu_{N,Q}(0)=0\) by
Lemma~\ref{lem:n-point-odd-moment-cancellation}.  Every summand vanishes.
\end{proof}

\begin{remark}[Even bump hypothesis is necessary]
\label{rem:n-point-even-bump-necessary}
The even-bump hypothesis in
Lemma~\ref{lem:n-point-convolution-parity} is load-bearing: if
\(\widehat\phi\) carries an odd component, the odd part of the product
no longer vanishes at any order below \(2N-1\).  An explicit numerical
witness for \(N=4\) appears in the simulation under
``scope check: non-even bump breaks parity vanishing.''
\end{remark}

\subsection{Algebraic-difference identity}
\label{subsec:n-point-difference-identity}

\begin{lemma}[Closed form for \(A_N(t)-A_N(t^{-1})\)]
\label{lem:A-N-difference}
With \(A_N(t):=\sum_{j=1}^N a_j^{(N)}\, t^j\),
\begin{equation}
\label{eq:A-N-difference}
        A_N(t) - A_N(t^{-1})
        = \frac{(-1)^N\,N}{\binom{2N-2}{N-1}}\,
          t^{-N}\,(1-t)^{2N-1}\,(1+t).
\end{equation}
\end{lemma}

\begin{proof}
Set \(F_N(t):=\sum_{j=-N}^N(-1)^j\binom{2N}{N+j}\,t^j\).  By the
binomial theorem, \(F_N(t)=(-1)^N t^{-N}(1-t)^{2N}\).  Differentiating
once and multiplying by \(t\),
\[
        t F_N'(t) = -N(-1)^N t^{-N}(1-t)^{2N-1}(1+t).
\]
On the other hand, the symmetry
\((-1)^{-j}\binom{2N}{N-j}=(-1)^j\binom{2N}{N+j}\) gives
\(tF_N'(t)=\sum_{j=1}^N j(-1)^j\binom{2N}{N+j}(t^j - t^{-j})\).
Multiplying by \(-1/\binom{2N-2}{N-1}\) and re-introducing the
\((-1)^{j+1}\) prefactor of \eqref{eq:n-point-coefficients} yields
\eqref{eq:A-N-difference}.
\end{proof}

\begin{corollary}[First surviving odd derivative]
\label{cor:n-point-first-surviving}
For every \(N\ge 1\) and every \(Q>0\),
\begin{equation}
\label{eq:mu-hat-first-surviving}
        \partial_z^{\,2N-1}\widehat\mu_{N,Q}(0)
        = \frac{(-1)^{N+1}(2N-1)\,N!\,(N-1)!}{Q^{2(2N-1)}}.
\end{equation}
\end{corollary}

\begin{proof}
Substituting \(t=e^{z/Q^2}\) in \eqref{eq:A-N-difference}, the
left-hand side equals
\(\widehat\mu_{N,Q}(z)-\widehat\mu_{N,Q}(-z)\).  Near \(z=0\),
\[
        1-e^{z/Q^2}=-\frac{z}{Q^2}+O(z^2),\quad
        e^{-Nz/Q^2}=1+O(z),\quad
        1+e^{z/Q^2}=2+O(z).
\]
Therefore
\[
        \widehat\mu_{N,Q}(z)-\widehat\mu_{N,Q}(-z)
        = \frac{(-1)^N\,N}{\binom{2N-2}{N-1}}\!
          \left(-\frac{z}{Q^2}\right)^{\!2N-1}\!\!\!\cdot 2
          + O(z^{2N})
        = \frac{2\,(-1)^{N+1}\,N\,z^{2N-1}}
               {\binom{2N-2}{N-1}\,Q^{2(2N-1)}}+O(z^{2N}).
\]
The \((2N-1)\)-fold derivative at \(z=0\) of the left-hand side is
\(2\,\partial_z^{2N-1}\widehat\mu_{N,Q}(0)\); on the right-hand side it
extracts \((2N-1)!\) times the leading coefficient.  Equating and using
\(\binom{2N-2}{N-1}=(2N-2)!/((N-1)!)^2\) together with
\(N(2N-1)!/(2N-2)!=N(2N-1)\) yields
\eqref{eq:mu-hat-first-surviving}.
\end{proof}

\begin{remark}[Where the order-\(2N-1\) vanishing comes from]
\label{rem:n-point-2N-minus-1-source}
The factor \((1-t)^{2N-1}\) in \eqref{eq:A-N-difference} is the
algebraic source of the order-\(2N-1\) odd-jet vanishing of
\(\widehat\mu_{N,Q}\): substituting \(t=e^{z/Q^2}\),
\(\widehat\mu_{N,Q}(z) - \widehat\mu_{N,Q}(-z)\) inherits a zero of
order exactly \(2N-1\) at \(z=0\), with leading coefficient given by
Corollary~\ref{cor:n-point-first-surviving}.
\end{remark}


\begin{remark}[Scope of \S\ref{sec:n-point-odd-moment-cancellation}]
\label{rem:n-point-scope}
This section supplies parity machinery only.  No analytic comparison,
finite-scale block estimate, source-energy transfer, or two-/four-point
exponent estimate appears here.  The downstream consumers of
Lemma~\ref{lem:n-point-coefficients} and
Lemma~\ref{lem:n-point-odd-moment-cancellation} must separately prove
that the scalar quantity to which \(\mu_{N,Q}\) is applied has a Taylor
expansion with the required even-bump/parity structure and with
remainders small enough for the relevant exponent comparison.  In
particular, this section does not prove the two-point or four-point
suppression hypotheses; it only provides the finite algebra used by those
later arguments.
\end{remark}
